\chapter{Figuras, Tabelas e Código}

\section{Figuras}

A \autoref{fig:exemplo} mostra como inserir uma figura. Guarde as imagens na
pasta \texttt{Imagens/}.

\begin{figure}[htbp]
  \centering
  \includegraphics[width=0.5\linewidth]{example-image}
  \caption{Exemplo de figura (a legenda fica por baixo)}
  \label{fig:exemplo}
  \fonte{Elaboração própria.}
\end{figure}

\section{Tabelas}

Nas tabelas a legenda fica por cima (\autoref{tab:exemplo}).

\begin{table}[htbp]
  \centering
  \caption{Exemplo de tabela}
  \label{tab:exemplo}
  \begin{tabular}{lcc}
    \toprule
    Método & Precisão (\%) & Tempo (s) \\
    \midrule
    A & 91,2 & 1,4 \\
    B & 94,7 & 2,9 \\
    \bottomrule
  \end{tabular}
  \fonte{Adaptado de \textcite{Resnick2009}.}
\end{table}

\section{Equações e código}

\begin{equation}
  E = mc^2
  \label{eq:exemplo}
\end{equation}

\begin{lstlisting}[language=Python, caption={Exemplo de código}, label={lst:exemplo}]
def media(valores):
    # Calcula a média (comentários com acentuação são suportados)
    return sum(valores) / len(valores)
\end{lstlisting}
